\section*{Capítulo \ref{ch:b1:ocean-circulation} --- La circulación del océano}
% values: calc/ch29.py

\begin{solution}{exo:b1:ocean-circulation:1}
(a)~A \qty{4}{m/s} (primera fórmula): $1000\,C_D = 0.29 + 0.775 + 0.481 = 1.546$,
$T = 1.3 \times \num{1.546e-3} \times 16 = \qty{0.0322}{Pa}$. A \qty{7}{m/s}
(segunda fórmula): $1000\,C_D = 0.60 + 0.497 = 1.097$, $T = \qty{0.0699}{Pa}$. A
\qty{15}{m/s}: $1000\,C_D = 1.665$, $T = \qty{0.487}{Pa}$. (b)~$0.487/0.0699 =
6.97$, frente a $(15/7)^2 = 4.59$: el coeficiente de arrastre mismo crece con el viento, porque un viento más fuerte levanta olas más rugosas que agarran mejor el aire.
\end{solution}

\begin{solution}{exo:b1:ocean-circulation:2}
(a)~Hemisferio norte, \ang{90} a la derecha de un viento que sopla hacia el este: hacia el sur. (b)~El viento sopla hacia el suroeste; a su derecha está el noroeste. (c)~Hemisferio sur, \ang{90} a la izquierda de un viento que sopla hacia el este: hacia el norte. (d)~El viento sopla hacia el norte; a su izquierda está el oeste, que es mar adentro en una costa occidental. En el caso (d) el transporte apunta mar adentro: el agua superficial se reemplaza desde abajo, y agua fría aflora a lo largo de la costa.
\end{solution}

\begin{solution}{exo:b1:ocean-circulation:3}
(a)~$31.7 + 3.5 - 17.9 = \qty{17.3}{\sverdrup}$, frente a
\qty{17.1}{\sverdrup} publicados: las medias de las tres partes suman el vuelco medio dentro del redondeo y de las incertidumbres ($\pm\qty{3.2}{\sverdrup}$).
(b)~Por encima de \qty{1100}{m}: $16.6 + 0.4 = \qty{17.0}{\sverdrup}$ hacia el norte; por debajo: $12.0 + 5.9 = \qty{17.9}{\sverdrup}$ hacia el sur. Los dos casi se equilibran, como deben hacerlo cuando no se acumula agua al norte de la sección; la pequeña diferencia cae dentro de las incertidumbres. (c)~$\num{17.1e6} \times \num{3.156e7} = \qty{5.40e14}{m^3}$, es decir
\qty{5.40e5}{km^3} al año; $\num{1.3324e9}/(\num{5.40e5}) \approx 2470$ años.
\end{solution}

\begin{solution}{exo:b1:ocean-circulation:4}
$D_E = 7.6 \times 10/\sqrt{\sin\varphi}$ y $V_0 = 0.0127 \times
10/\sqrt{\sin\varphi}$. A \ang{15}: \qty{149}{m} y \qty{0.250}{m/s} (\qty{2.5}{\percent} del viento); a \ang{45}: \qty{90}{m} y \qty{0.151}{m/s} (\qty{1.5}{\percent}); a \ang{60}: \qty{82}{m} y \qty{0.136}{m/s} (\qty{1.4}{\percent}). Cerca del ecuador $\sin\varphi \to 0$ y ambas fórmulas divergen; más fundamentalmente $f \to 0$, así que la fuerza de Coriolis necesaria para equilibrar la fricción se anula, y el tiempo $1/f$ que necesita el equilibrio para establecerse se vuelve mayor que el tiempo que los vientos se mantienen constantes. Ekman restringió su fórmula a $|\varphi| \geq
\ang{10}$.
\end{solution}

\begin{solution}{exo:b1:ocean-circulation:5}
(a)~$f = 2\Omega\sin\ang{11} = \qty{2.783e-5}{s^{-1}}$; $M = 0.07/(\num{2.783e-5}) =
\qty{2515}{kg.m^{-1}.s^{-1}}$, a \ang{90} a la derecha de un viento que sopla hacia el oeste: hacia el norte. (b)~$Q = (2515/1027) \times \num{5e6} = \qty{1.22e7}{m^3/s}$,
\qty{12.2}{\sverdrup}. (c)~La estimación cae dentro de la incertidumbre de la medida directa ($12.0 \pm 5.5$) y cerca del valor a partir de los vientos de largo plazo ($13.5 \pm 0.3$): el teorema explica el transporte observado.
\end{solution}

\begin{solution}{exo:b1:ocean-circulation:6}
(a)~$f(\ang{2}) = \qty{5.09e-6}{s^{-1}}$; $M = 0.05/(\num{5.09e-6}) =
\qty{9824}{kg.m^{-1}.s^{-1}}$, \qty{9.57}{m^2/s}: a \ang{2} N hacia el norte (a la derecha de un viento del oeste), a \ang{2}~S hacia el sur (a la izquierda). (b)~Sobre una franja de \qty{1}{m} de longitud en longitud, $2 \times 9.57 = \qty{19.1}{m^3/s}$ salen a través de las dos \omterm{def:b1:spherical-astronomy:geographic}{latitudes}, separadas \qty{444.8}{km}: $w = 19.1/(\num{4.448e5}) = \qty{4.30e-5}{m/s}$,
unos \qty{3.7}{m} al día. (c)~En el ecuador $f = 0$ y $|\vec M| = T/|f|$ es infinito: el equilibrio de Ekman estacionario no existe allí, y las cifras anteriores, muy sensibles a la \omterm{def:b1:spherical-astronomy:geographic}{latitud} escogida, son solo un orden de magnitud. Lo que sobrevive es el signo: a ambos lados el transporte apunta lejos del ecuador, así que la superficie diverge y agua fría asciende a lo largo de él.
\end{solution}

\begin{solution}{exo:b1:ocean-circulation:7}
(a)~$f(\ang{15}) = \qty{3.775e-5}{s^{-1}}$: $M = 0.06/(\num{3.775e-5}) =
\qty{1590}{kg.m^{-1}.s^{-1}}$ hacia el norte; $f(\ang{40}) =
\qty{9.375e-5}{s^{-1}}$: $M = \qty{1280}{kg.m^{-1}.s^{-1}}$ hacia el sur.
(b)~Entre las dos \omterm{def:b1:spherical-astronomy:geographic}{latitudes}, separadas \qty{2780}{km}, $\partial M_y/\partial y =
-(1590 + 1280)/(\num{2.78e6}) = \qty{-1.03e-3}{kg.m^{-2}.s^{-1}}$: por \cref{prop:b1:ocean-circulation:pumping}, $w_E = \qty{-1.005e-6}{m/s}$, una velocidad hacia abajo de \qty{31.7}{m} al año. (c)~$\num{1.005e-6} \times \num{2.78e6} \times
\num{5e6} = \qty{1.40e7}{m^3/s}$, \qty{14.0}{\sverdrup} bombeados hacia abajo.
\end{solution}

\begin{solution}{exo:b1:ocean-circulation:8}
(a)~$H = 1027 \times 3992 \times \num{17e6} \times 12 = \qty{8.36e14}{W}$,
\qty{0.84}{PW}. (b)~Es el \qty{70}{\percent} del \qty{1.2}{PW} medido: la rama superior real es más cálida cerca de la superficie que la media de \qty{15}{\celsius}, y el \omterm{def:b1:ocean-circulation:gyre}{giro} también transporta calor. (c)~$1027 \times 3992 \times \num{17e6} \times 3 =
\qty{2.1e14}{W}$, \qty{0.21}{PW}, una cuarta parte de (a). Solo el agua que se enfría antes de regresar transporta mucho calor: el transporte de calor del Atlántico depende de la \omterm{def:b1:ocean-circulation:moc}{formación de agua profunda}.
\end{solution}

\begin{solution}{exo:b1:ocean-circulation:9}
Sean $x$, $y$ y $1 - x - y$ las fracciones de AAIW, NADW y AABW. Balances de sal y calor, medidos desde AABW:
\[
  -0.52\,x + 0.33\,y = 34.731 - 34.67 = 0.061,\qquad
  2.9\,x + 3.6\,y = 2.78 - 0.4 = 2.38 .
\]
El determinante es $-0.52 \times 3.6 - 0.33 \times 2.9 = -2.829$, y la regla de Cramer da $x = (0.061 \times 3.6 - 0.33 \times 2.38)/(-2.829) = 0.20$ y $y = (-0.52
\times 2.38 - 2.9 \times 0.061)/(-2.829) = 0.50$: la muestra es \qty{20}{\percent} AAIW, \qty{50}{\percent} NADW y \qty{30}{\percent} AABW, un punto dentro del triángulo de las tres aguas.
\end{solution}

\begin{solution}{exo:b1:ocean-circulation:10}
(a)~Con $W = u + iv = V_0e^{i\pi/4}e^{(1+i)az}$, cuyas partes real e imaginaria son
las fórmulas de la proposición, $W'' = (1+i)^2a^2W = 2ia^2W$. Las dos ecuaciones
se combinan en $A_z(u'' + iv'') = -fv + ifu = if(u + iv)$, es decir $A_zW'' = ifW$,
que se cumple porque $2ia^2A_z = if$ para $a^2 = f/(2A_z)$. (b)~$W'(0) =
(1+i)aV_0e^{i\pi/4} = \sqrt2\,aV_0\,e^{i\pi/2} = i\sqrt2\,aV_0$: la tensión de superficie
$\rho A_zW'(0)$ apunta hacia el norte, con el tamaño $\rho A_z\sqrt2\,aV_0 = \rho A_z\sqrt2
\sqrt{f/(2A_z)}\;T/(\rho\sqrt{fA_z}) = T$. (c)~$\int_{-\infty}^0 W\,\mathrm dz =
V_0e^{i\pi/4}/((1+i)a) = V_0/(\sqrt2\,a)$, un número real: el transporte no tiene
parte hacia el norte, y su parte hacia el este es $\rho V_0/(\sqrt2\,a) =
T/\sqrt{fA_z}\cdot\sqrt{A_z/f} = T/f$. Este es el teorema: tamaño $T/f$, \ang{90} a la
derecha de un viento que sopla hacia el norte.
\end{solution}

\begin{solution}{exo:b1:ocean-circulation:11}
(a)~Un paso hacia el norte $\mathrm dy = a_\oplus\,\mathrm d\varphi$ cambia $f = 2\Omega
\sin\varphi$ en $2\Omega\cos\varphi\,\mathrm d\varphi$, así que $\beta = 2\Omega
\cos\varphi/a_\oplus = 2 \times \num{7.292e-5} \times \cos\ang{30}/(\num{6.371e6}) =
\qty{1.98e-11}{m^{-1}.s^{-1}}$. (b)~$R/\beta = \num{e-6}/(\num{1.98e-11}) =
\qty{5.04e4}{m}$, unos \qty{50}{km}. (c)~Interior: $\num{30e6}/(\num{5e6} \times
1000) = \qty{6.0e-3}{m/s}$, \qty{6}{mm/s}; corriente de frontera: $\num{30e6}/(\num{5.04e4}
\times 1000) = \qty{0.60}{m/s}$, cien veces más rápida, el cociente de los anchos. La corriente del Golfo, a \qty{1.8}{m/s}, es aún más rápida: la corriente real es más estrecha que esta estimación, su flujo se concentra cerca de la superficie, y su recirculación se suma a su transporte.
\end{solution}

\begin{solution}{exo:b1:ocean-circulation:12}
(a)~Donde el mar es más profundo que \qty{1000}{m}, la columna por encima de \qty{1000}{m} contiene como mucho \qty{1}{km} de agua, y en otras partes menos: el volumen por encima de \qty{1000}{m} es como mucho $\num{361.84e6} \times 1 = \qty{3.618e8}{km^3}$, así que el volumen por debajo es al menos $\num{1.3324e9} - \num{3.618e8} = \qty{9.71e8}{km^3}$.
(b)~$\num{26e6} \times \num{3.156e7} = \qty{8.21e14}{m^3}$ al año,
\qty{8.21e5}{km^3}; $\num{9.71e8}/(\num{8.21e5}) \approx 1200$ años, del mismo orden que los mil años de la cinta transportadora. (c)~$(1/2)^{25/12.3} = 0.244$, cerca de una cuarta parte;
$\num{e-3} \times 25 \times \num{3.156e7} = \qty{7.9e5}{m}$, unos \qty{790}{km}. Tras 25 años el tritio se ha extendido no más lejos de lo que lo llevaría una deriva lenta, mientras que queda una cuarta parte de él para detectar.
\end{solution}

\begin{solution}{pb:b1:ocean-circulation:1}
\begin{enumerate}
  \item Un viento del norte sopla hacia el sur, a lo largo de la costa.
  \item $1000\,C_D = 0.60 + 0.071 \times 8 = 1.168$: $C_D = \num{1.168e-3}$.
  \item $T = 1.3 \times \num{1.168e-3} \times 64 = \qty{0.0972}{Pa}$.
  \item $C_D = \num{1.452e-3}$ y $T = \qty{0.272}{Pa}$: la tensión crece en $2.80$,
        más que $(12/8)^2 = 2.25$, porque $C_D$ mismo crece con el viento.
  \item $\Omega = \qty{7.292e-5}{rad.s^{-1}}$; $f = 2\Omega\sin\ang{38} =
        \qty{8.979e-5}{s^{-1}}$.
  \item $2\pi/f = \qty{7.00e4}{s}$, \qty{19.4}{h}. La fuerza de Coriolis hace dar a una corriente una vuelta completa en este tiempo; antes de que el viento haya soplado durante
        cerca de ese tiempo, el agua oscila en lugar de asentarse en el equilibrio estacionario del teorema.
  \item \ang{90} a la derecha de un viento que sopla hacia el sur: hacia el oeste, el mar abierto.
  \item $M = T/f = 0.0972/(\num{8.979e-5}) = \qty{1082}{kg.m^{-1}.s^{-1}}$.
  \item $1082/1027 = \qty{1.054}{m^2/s}$ por cada metro de costa.
  \item $D_E = 7.6 \times 8/\sqrt{\sin\ang{38}} = \qty{77.5}{m}$.
  \item $1.054/77.5 = \qty{0.0136}{m/s}$, \qty{1.36}{cm/s}.
  \item $V_0 = 0.0127 \times 8/\sqrt{\sin\ang{38}} = \qty{0.129}{m/s}$, \ang{45} a la derecha del viento: hacia el suroeste. Las boyas dan cerca de la mitad del valor de Ekman, unos \qty{0.07}{m/s}.
  \item El agua que sale de la costa en la \omterm{def:b1:ocean-circulation:ekman}{capa de Ekman} no puede reemplazarse a través de la costa, ni desde arriba: se reemplaza con agua que sube desde debajo de la capa cerca de la costa, como en la sección transversal dibujada en \cref{sec:b1:ocean-circulation:upwelling}; las isotermas se curvan hacia arriba hacia la orilla.
  \item $1.054 \times 1000 \times 86400 = \qty{9.11e7}{m^3}$ al día por kilómetro de costa.
  \item $w = 1.054/(\num{2e4}) = \qty{5.27e-5}{m/s}$, \qty{4.55}{m} al día.
  \item $77.5/4.55 = 17$ días.
  \item El aire se enfría desde abajo por el agua fría: la capa cercana a la superficie se vuelve fría y estable, la humedad puede condensarse en niebla y estratos bajos, y sin convección hay poca lluvia en la costa.
  \item El agua afluida es rica en nutrientes; estos alimentan al fitoplancton, que alimenta al zooplancton y a los peces.
  \item Un viento que sopla hacia el norte impulsa un transporte a \ang{90} a su derecha, hacia el este: el agua superficial se acumula contra la costa y se hunde (\omterm{def:b1:ocean-circulation:upwelling}{hundimiento}); el \omterm{def:b1:ocean-circulation:upwelling}{afloramiento} se detiene y el agua superficial más cálida regresa a la costa.
  \item $f(\ang{20}) = \qty{4.988e-5}{s^{-1}}$: $M = \qty{1948}{kg.m^{-1}.s^{-1}}$,
        $\sin\ang{38}/\sin\ang{20} = 1.80$ veces más.
  \item En el hemisferio sur el transporte está a \ang{90} a la izquierda del viento. Para empujarlo hacia el oeste, mar adentro, el viento debe soplar hacia el norte, desde el sur: en ambos casos el viento sopla a lo largo de la costa hacia el ecuador.
  \item $1000\,C_D = 0.29 + 0.62 + 0.308 = 1.218$; $T = \qty{0.0396}{Pa}$; $M =
        \qty{441}{kg.m^{-1}.s^{-1}}$; el cociente es $0.407$ en ambos casos.
  \item Por metro de costa, $M c_p\,\Delta T = 1082 \times 3992 \times 6 =
        \qty{2.59e7}{W}$, \qty{25.9}{MW}; a lo largo de \qty{1000}{km}, \qty{2.59e13}{W},
        \qty{25.9}{TW}.
  \item $\num{31.7e6}/1.054 = \qty{3.01e7}{m}$, unos \qty{30000}{km}: el \omterm{def:b1:ocean-circulation:upwelling}{afloramiento costero} mueve mucha menos agua que una \omterm{def:b1:ocean-circulation:wbc}{corriente de frontera occidental}.
  \item $Q = 1.054 \times \num{e6} = \qty{1.05e6}{m^3/s}$: el transporte de \omterm{def:b1:ocean-circulation:upwelling}{afloramiento} a lo largo de \qty{1000}{km} de costa es de \qty{1.05}{\sverdrup}.
\end{enumerate}
\end{solution}
